\scanpage{122}
% line head | text | 71
\mspage{71}
% line 1–3 | text | déduites de cette formule
déduites de cette formule ne sont autres que \deleted{celles}
les suivantes : $u,t$ sont donnés comme objets de
$\mathord{}_{u'}\backslash^{U}$, $T_{/t'}$, i.e. avec des flèches
% line 4 | display | alpha and beta
\[
u'\xrightarrow{\alpha}f_U(u),\qquad
\deleted{\unsure{f}}f_T(t)\mathrel{\overset{\beta}{\unsure{-|\to}}}t'
\]
% line 5 | text | d'où
d'où
% line 6 | diagram | Hom_X via f, then via alpha and beta
\diagram{\begin{tikzcd}[ampersand replacement=\&,column sep=small]
\operatorname{Hom}_X(u,t)
\arrow[r,"{\text{via }\deleted{\unsure{g}}f=f_X:X\to X'}"]
\& \operatorname{Hom}_{X'}(\underset{f_U(u)}{\underline{f(u)}},\underset{f_T(t)}{\underline{f(t)}})
\arrow[r,"{\text{via }\alpha{,}\beta}"]
\& \operatorname{Hom}_{X'}(u',t').
\end{tikzcd}}

% line 7–9 | text | Notons pour mémoire
Notons pour mémoire que pour tout
objet $Z$ de $\mathrm{Cat}$, et tout préfaisceau $G$
sur $Z$, on a la formule
% line 10 | display | lim G and pi nought
\[
\underset{Z}{\lim_{\rightarrow}}G\simeq\pi_0(Z_{/G}),
\]
% line 11–23 | text | où Z/G s'interprète
où $Z_{/G}$ s'interprète en terme\unsure{s} de l'immersion
canonique $Z\longrightarrow Z^\wedge$ ($\unsure{C|G}\in\operatorname{Ob}Z^\wedge$).
C'est donc une catégorie \deleted{\unsure{co}}fibrée sur $Z$ à fibres
discrètes, \unsure{\{}et on a \inserted{en fait} un \deleted{équivalence}\inserted{isom.} entre la
\inserted{ss} catégorie \inserted{pleine} de $\unsure{\mathrm{Kat}|\mathrm{Cat}}_{/Z}$ formée \unsure{de|des} telles catégories
fibrées à fibres discrètes, avec $Z^\wedge\unsure{\}|\big)}$.
Dans ce dictionnaire parfait entre préfaisceaux
d'une\unsure{.}, et cat. fibrées à fibres discrètes,
le foncteur $\lim_{\rightarrow}$ sur $Z^\wedge$ correspond simplement
au \inserted{foncteur} $\pi_0$ \deleted{\unsure{dans}} de la catégorie des objets au
dessus de $Z$ dans $\mathrm{Cat}$, lui-même induit
par le foncteur (indépendant de $Z$)
% line 24 | display | Cat to (Ens), pi nought
\[
\mathrm{Cat}\xrightarrow{\pi_0}(\mathrm{Ens}),
\]
% archive sheet number: 122
